\section{A local compute-efficiency threshold}
The break-even identity is retrospective: it asks whether the wide phase has already created enough discovery advantage to repay its historical compute. For an online controller we also need a prospective question: starting from the same current predictor, is the next infinitesimal unit of compute better spent in the wide parameter space or in a reduced subspace? In the ideal spectral model this comparison is exact and produces a dimensionless threshold.

Let $g=\nabla_\theta\R(\theta)$ at the current checkpoint and let $Q$ be an orthogonal projector onto the parameter directions retained by a candidate contraction. Consider two continuous-time updates starting from the same parameter point,
\[
\dot\theta_W=-g(\theta_W),
\qquad
\dot\theta_Q=-Q g(\theta_Q).
\]
Suppose one unit of gradient-flow time costs $C_W$ for the full update and $C_Q$ for the projected update, with $0<C_Q\le C_W$. Parameterize both trajectories by cumulative compute $b$, so that an increment $db$ corresponds locally to $db/C_W$ units of wide flow or $db/C_Q$ units of projected flow.

\begin{theorem}[Exact infinitesimal value per unit compute]\label{thm:local-compute-efficiency}
At the common checkpoint,
\[
\left.\frac{d}{db}\R(\theta_W(b))\right|_{b=0}
=-\frac{\norm{g}^2}{C_W},
\qquad
\left.\frac{d}{db}\R(\theta_Q(b))\right|_{b=0}
=-\frac{\norm{Qg}^2}{C_Q}.
\]
Therefore the reduced update has strictly larger instantaneous risk decrease per unit compute if and only if
\[
\boxed{
\frac{\norm{Qg}^2}{C_Q}
>
\frac{\norm{g}^2}{C_W}.}
\]
Writing
\[
A_S:=\norm{Qg}^2,
\qquad
A_D:=\norm{(I-Q)g}^2,
\qquad
A=A_S+A_D,
\]
this condition is equivalently
\[
\boxed{
A_D
<
\left(\frac{C_W}{C_Q}-1\right)A_S,}
\]
or, whenever $A>0$,
\[
\boxed{
\frac{A_D}{A}
<
1-\frac{C_Q}{C_W}.}
\]
Thus the fraction of instantaneous gradient-descent power discarded by contraction must be smaller than the fraction of per-update compute saved.
\end{theorem}

\begin{proof}
Along ordinary gradient flow,
\[
\frac{d}{ds}\R(\theta_W(s))
=
\ip{g}{-g}
=-\norm{g}^2.
\]
Since $s=b/C_W$, the chain rule gives the first derivative. Along the projected flow,
\[
\frac{d}{ds}\R(\theta_Q(s))
=
\ip{g}{-Qg}
=-\norm{Qg}^2,
\]
because $Q$ is an orthogonal projector; reparameterizing by $s=b/C_Q$ gives the second derivative. Since $Q$ and $I-Q$ have orthogonal ranges,
$\norm{g}^2=A_S+A_D$. The reduced update is locally more compute-efficient exactly when
$A_S/C_Q>(A_S+A_D)/C_W$. Rearranging gives the second boxed inequality. Dividing by $A$ gives the last one.
\end{proof}

\begin{corollary}[Target-weighted spectral form]\label{cor:local-compute-spectral}
Let $J=U\Sigma V^\ast$, let $\mu_j$ be the positive eigenvalues of $JJ^\ast$, and let $a_j=\ip{e}{u_j}$. If $Q$ retains the right singular directions indexed by $S$ and drops $D=S^c$, then
\[
A_S=\sum_{j\in S}\mu_j a_j^2,
\qquad
A_D=\sum_{j\in D}\mu_j a_j^2.
\]
Consequently the candidate is instantaneously more efficient per unit compute if and only if
\[
\boxed{
\frac{\sum_{j\in D}\mu_j a_j^2}
{\sum_j\mu_j a_j^2}
<
1-\frac{C_Q}{C_W}.}
\]
\end{corollary}

\begin{proof}
The modal gradient identity gives
$g=J^\ast e=\sum_j\sqrt{\mu_j}a_jv_j$.
Orthogonality of the $v_j$ yields the two energy sums. Apply Theorem~\ref{thm:local-compute-efficiency}.
\end{proof}

\begin{remark}[Why this is not yet a pruning theorem]
The inequality above answers only the next-infinitesimal-compute question. It does not control the structural jump produced by a physical contraction, the approximation floor created by permanently deleting target mass, or the possibility that currently weak directions later become target-relevant because the Jacobian rotates. Those are controlled by the structural-transfer, finite-horizon, and future-discovery certificates developed elsewhere in this note. The useful interpretation is therefore not ``prune whenever the inequality holds,'' but rather: among candidates that are already statistically and geometrically safe, reject any contraction that cannot even beat the wide model in instantaneous risk decrease per unit compute.
\end{remark}

\begin{remark}[A dimensionless onset statistic]
Define the omitted descent-power fraction
\[
\omega_D(t):=
\frac{\norm{(I-Q_t)\nabla\R(\theta_t)}^2}
{\norm{\nabla\R(\theta_t)}^2}
\]
and the structural compute-saving fraction
\[
s_C(t):=1-\frac{C_{Q_t}}{C_W}.
\]
The local compute gate is simply $\omega_D(t)<s_C(t)$. Both sides are dimensionless. The left side is target-weighted prediction geometry; the right side is a property of the actual structural implementation. This is the cleanest local bridge between the spectral theory and physical compute.
\end{remark}

\begin{corollary}[First-order break-even over a small compute budget]\label{cor:local-break-even}
Assume the two trajectories above are differentiable in compute at $b=0$. For a small equal-compute budget $B>0$,
\[
\R(\theta_Q(B))-\R(\theta_W(B))
=
B\left(
\frac{\norm{g}^2}{C_W}
-
\frac{\norm{Qg}^2}{C_Q}
\right)+o(B).
\]
Hence the sign in Theorem~\ref{thm:local-compute-efficiency} is the exact first-order sign of the equal-compute risk difference.
\end{corollary}

\begin{proof}
Apply the first-order Taylor expansion in the common compute parameter $b$ to both risk trajectories and subtract.
\end{proof}
